\documentclass[11pt]{article}

\usepackage[margin=1in]{geometry}
\usepackage{amsmath, amssymb, amsthm, mathtools, bm}
\usepackage{graphicx}
\usepackage{algorithm}
\usepackage{float}
\usepackage{algpseudocode}
\usepackage{hyperref}
\usepackage{listings}
\usepackage{xcolor}
\usepackage{booktabs}
\usepackage{multirow}

\lstdefinestyle{python}{
  language=Python,
  basicstyle=\ttfamily\small,
  keywordstyle=\bfseries\color{blue},
  commentstyle=\itshape\color{green!60!black},
  stringstyle=\color{red},
  showstringspaces=false,
  frame=single,
  breaklines=true,
  numbers=left,
  numberstyle=\tiny\color{gray}
}

\title{TurtleBot Interceptor: Complete System Implementation\\
\large From Simulation to Hardware Deployment}
\author{Tim, Gabe, Thomas, Ryan, Kush}
\date{}

\begin{document}
\maketitle

\tableofcontents
\newpage

\section{Introduction}

This document provides a comprehensive technical description of the TurtleBot Interceptor system, covering algorithm design, implementation details, simulation architecture, and hardware deployment. The system successfully intercepts a moving target robot in unknown, cluttered environments using Monte Carlo Localization (MCL), Kalman Filtering, and Model Predictive Control (MPC).

\subsection{System Overview}

The system consists of four main layers:

\begin{enumerate}
    \item \textbf{State Estimation Layer}: Extended Kalman Filter (EKF) fusing IMU, magnetometer, and wheel encoders
    \item \textbf{Mapping Layer}: Cartographer SLAM + Fast Local Grid for persistent environment mapping
    \item \textbf{Target Tracking Layer}: Multi-sensor fusion (camera + LIDAR) with robust filtering
    \item \textbf{Control Layer}: Uncertainty-aware MPC with obstacle avoidance
\end{enumerate}

\textbf{Localization Approach:} The system uses \textbf{EKF for robot localization} (not MCL) in hardware deployment. MCL is available for simulation but is commented out in the main launch file.

\section{State Estimation Architecture}

\subsection{Dual Estimation Design}

We employ a dual estimation architecture that separates high-rate state estimation from mapping:

\subsubsection{Extended Kalman Filter (EKF)}

The EKF provides precise, high-rate robot state estimation at 100 Hz, essential for smooth control.

\textbf{State Vector:}
\[
\mathbf{x} = \begin{bmatrix} x \\ y \\ \theta \\ v_x \\ v_y \\ \omega \end{bmatrix}
\]

\textbf{Process Model:}
\begin{align}
x_{k+1} &= x_k + \Delta t (v_{x,k} \cos\theta_k - v_{y,k} \sin\theta_k) + w_{x,k} \\
y_{k+1} &= y_k + \Delta t (v_{x,k} \sin\theta_k + v_{y,k} \cos\theta_k) + w_{y,k} \\
\theta_{k+1} &= \theta_k + \Delta t \cdot \omega_k + w_{\theta,k} \\
v_{x,k+1} &= v_{x,k} + w_{v_x,k} \\
v_{y,k+1} &= v_{y,k} + w_{v_y,k} \\
\omega_{k+1} &= \omega_k + w_{\omega,k}
\end{align}

where $\mathbf{w}_k \sim \mathcal{N}(0, \mathbf{Q})$ is process noise.

\textbf{Measurement Models:}

\begin{itemize}
    \item \textbf{IMU Gyroscope}: Direct measurement of $\omega$
    \item \textbf{IMU Accelerometer}: Gravity-compensated acceleration
    \[
    \mathbf{a}_{\text{world}} = \mathbf{R}(\theta) \mathbf{a}_{\text{IMU}} - \mathbf{g}
    \]
    where $\mathbf{R}(\theta)$ is the rotation matrix and $\mathbf{g} = [0, 0, 9.81]^T$.
    \item \textbf{Magnetometer}: Absolute heading measurement
    \[
    \theta_{\text{mag}} = \atan2(m_y, m_x) - \theta_{\text{offset}}
    \]
    where $\theta_{\text{offset}}$ is calibrated at startup.
    \item \textbf{Wheel Encoders}: Direct velocity measurements
\end{itemize}

\textbf{Process Noise Covariance:}
\[
\mathbf{Q} = \text{diag}([0.001, 0.001, 0.01, 0.05, 0.05, 0.05])
\]

\textbf{Measurement Noise Covariances:}
\begin{align}
R_{\text{gyro}} &= 0.01 \\
R_{\text{accel}} &= 0.5 \\
R_{\text{mag}} &= 0.1 \\
R_{\text{odom}} &= 0.02
\end{align}

\subsubsection{Cartographer SLAM}

Cartographer provides robust, persistent mapping at 1-2 Hz update rate. It uses:
\begin{itemize}
    \item EKF odometry as a constraint
    \item LIDAR scan matching for fine alignment
    \item Sub-mapping for efficiency
    \item Loop closure for drift correction
\end{itemize}

\textbf{Map Resolution:} 2 cm (0.02 m)

\textbf{Output:} OccupancyGrid with values:
\begin{itemize}
    \item $-1$: Unknown
    \item $0$: Free space
    \item $100$: Occupied
\end{itemize}

\section{Target State Estimation}

\subsection{Multi-Source Fusion Architecture}

The target estimator combines two complementary sensing modalities:

\subsubsection{Vision-Based Detection}

\textbf{HSV Color Detection:}
\begin{itemize}
    \item Yellow cones (obstacles): HSV $[15, 80, 80]$ to $[35, 255, 255]$
    \item Blue cone (target): HSV $[100, 50, 50]$ to $[130, 255, 255]$
\end{itemize}

\textbf{Depth Estimation:}
Using known cone area and pixel count:
\[
\text{depth} = \sqrt{\frac{f_x f_y \cdot \text{CONE\_AREA}}{\text{pixel\_count}}}
\]

where:
\begin{itemize}
    \item Yellow cones: $\text{CONE\_AREA} = 0.0208$ m$^2$ (ground level)
    \item Blue target: $\text{TARGET\_CONE\_AREA} = 0.146$ m$^2$ (7$\times$ larger, elevated $\approx 10$ cm above ground)
\end{itemize}

\textbf{Target Cone Elevation:} The blue target cone is positioned approximately 10 cm above the ground ($z \approx 0.1$ m), distinguishing it from yellow obstacle cones which sit at ground level. This elevation is accounted for in the coordinate transform from camera frame to map frame.

\textbf{3D to Map Transform:}
\begin{enumerate}
    \item Camera frame $\rightarrow$ base\_link: $\mathbf{p}_{\text{base}} = \mathbf{T}_{\text{cam}\rightarrow\text{base}} \mathbf{p}_{\text{cam}}$
    \item base\_link $\rightarrow$ map: 
    \[
    \mathbf{p}_{\text{map}} = \mathbf{R}(\theta_{\text{robot}}) \mathbf{p}_{\text{base}} + \mathbf{p}_{\text{robot}}
    \]
    \item \textbf{Target elevation}: Blue target cone at $z \approx 0.1$ m (10 cm above ground), yellow cones at $z \approx 0$ m
\end{enumerate}

\subsubsection{LIDAR-Based Detection}

When both robots have map origins available, we extract direct robot positions and compute the transform using the Kabsch algorithm (SVD-based).

\textbf{Kabsch Algorithm:}
Given two sets of corresponding points $\mathbf{A} = \{\mathbf{a}_i\}$ and $\mathbf{B} = \{\mathbf{b}_i\}$:

\begin{enumerate}
    \item Compute centroids:
    \[
    \bar{\mathbf{a}} = \frac{1}{n}\sum_{i=1}^n \mathbf{a}_i, \quad \bar{\mathbf{b}} = \frac{1}{n}\sum_{i=1}^n \mathbf{b}_i
    \]
    \item Center the points:
    \[
    \mathbf{A}_{\text{cent}} = \mathbf{A} - \bar{\mathbf{a}}, \quad \mathbf{B}_{\text{cent}} = \mathbf{B} - \bar{\mathbf{b}}
    \]
    \item Compute cross-covariance:
    \[
    \mathbf{H} = \mathbf{A}_{\text{cent}}^T \mathbf{B}_{\text{cent}}
    \]
    \item SVD decomposition:
    \[
    \mathbf{H} = \mathbf{U} \mathbf{\Sigma} \mathbf{V}^T
    \]
    \item Rotation matrix:
    \[
    \mathbf{R} = \mathbf{V}^T \mathbf{U}^T
    \]
    If $\det(\mathbf{R}) < 0$, reflect: $\mathbf{R} = \mathbf{V}^T \mathbf{M} \mathbf{U}^T$ where $\mathbf{M} = \text{diag}([1, 1, -1])$
    \item Translation:
    \[
    \mathbf{t} = \bar{\mathbf{b}} - \mathbf{R} \bar{\mathbf{a}}
    \]
\end{enumerate}

\subsection{Robust Filtering Pipeline}

\subsubsection{Temporal Smoothing}

\textbf{Exponential Moving Average:}
\[
\mathbf{x}_{\text{filtered},k} = (1-\alpha) \mathbf{x}_{\text{filtered},k-1} + \alpha \mathbf{x}_{\text{raw},k}
\]

where $\alpha_{\text{position}} = 0.8$ (favor latest detection) and $\alpha_{\text{orientation}} = 0.3$.

\textbf{Median Filtering:}
\[
\mathbf{x}_{\text{robust}} = \text{median}(\{\mathbf{x}_i\}_{i=1}^{N})
\]

over the last $N=10$ detections.

\subsubsection{Outlier Rejection}

\textbf{RANSAC-like Algorithm:}
\begin{enumerate}
    \item Randomly sample $m$ point pairs (minimum for transform estimation)
    \item Compute transform from sample
    \item Count inliers (points within threshold distance)
    \item Repeat for $K$ iterations
    \item Use transform with most inliers
\end{enumerate}

\textbf{Parameters:}
\begin{itemize}
    \item Max outliers: 60\%
    \item Trials: 50
    \item Inlier threshold: 25 cm
\end{itemize}

\subsubsection{State Propagation When Target is Out of Frame}

When the target is not visible, we propagate the state forward using a constant-velocity model:

\textbf{State Vector:}
\[
\mathbf{x}_{\text{target}} = \begin{bmatrix} x \\ y \\ v_x \\ v_y \\ \theta \\ \omega \end{bmatrix}
\]

\textbf{Propagation Model:}
\begin{align}
x_{k+1} &= x_k + \Delta t \cdot v_{x,k} + w_{x,k} \\
y_{k+1} &= y_k + \Delta t \cdot v_{y,k} + w_{y,k} \\
v_{x,k+1} &= v_{x,k} + w_{v_x,k} \\
v_{y,k+1} &= v_{y,k} + w_{v_y,k} \\
\theta_{k+1} &= \theta_k + \Delta t \cdot \omega_k + w_{\theta,k} \\
\omega_{k+1} &= \omega_k + w_{\omega,k}
\end{align}

\textbf{Covariance Propagation:}
\[
\mathbf{P}_{k+1} = \mathbf{F}_k \mathbf{P}_k \mathbf{F}_k^T + \mathbf{Q}_k
\]

where $\mathbf{F}_k$ is the Jacobian of the propagation model:
\[
\mathbf{F}_k = \begin{bmatrix}
1 & 0 & \Delta t & 0 & 0 & 0 \\
0 & 1 & 0 & \Delta t & 0 & 0 \\
0 & 0 & 1 & 0 & 0 & 0 \\
0 & 0 & 0 & 1 & 0 & 0 \\
0 & 0 & 0 & 0 & 1 & \Delta t \\
0 & 0 & 0 & 0 & 0 & 1
\end{bmatrix}
\]

\textbf{Process Noise:}
\[
\mathbf{Q}_k = \text{diag}([\sigma_x^2, \sigma_y^2, \sigma_{v_x}^2, \sigma_{v_y}^2, \sigma_{\theta}^2, \sigma_{\omega}^2])
\]

where noise scales with time since last detection:
\[
\sigma_i^2 = \sigma_{i,\text{base}}^2 \cdot (1 + \beta \cdot \Delta t_{\text{lost}})
\]

with $\beta = 0.5$ to increase uncertainty over time.

\textbf{Confidence Decay:}
\[
\text{confidence}_{k+1} = \text{confidence}_k \cdot e^{-\lambda \Delta t}
\]

where $\lambda = 0.1$ s$^{-1}$ controls decay rate.

\section{Model Predictive Control}

\subsection{Unicycle Dynamics}

\textbf{State:}
\[
\mathbf{x} = \begin{bmatrix} p_x \\ p_y \\ \theta \\ v \end{bmatrix}
\]

\textbf{Control:}
\[
\mathbf{u} = \begin{bmatrix} a \\ \omega \end{bmatrix}
\]

\textbf{Discrete-Time Dynamics:}
\begin{align}
p_{x,k+1} &= p_{x,k} + \Delta t \cdot v_k \cos\theta_k \\
p_{y,k+1} &= p_{y,k} + \Delta t \cdot v_k \sin\theta_k \\
\theta_{k+1} &= \theta_k + \Delta t \cdot \omega_k \\
v_{k+1} &= v_k + \Delta t \cdot a_k
\end{align}

\subsection{Optimization Problem}

At each time step $t$, solve:
\begin{align}
\min_{\mathbf{u}_{0:N-1}} \quad & J = \sum_{k=0}^{N-1} \ell(\mathbf{x}_k, \mathbf{u}_k, \mathbf{x}_k^{\text{tgt}}) + \ell_f(\mathbf{x}_N, \mathbf{x}_N^{\text{tgt}}) \\
\text{s.t.} \quad & \mathbf{x}_{k+1} = f(\mathbf{x}_k, \mathbf{u}_k), \quad k = 0, \ldots, N-1 \\
& \mathbf{x}_0 = \hat{\mathbf{x}}_t^{\text{seek}} \\
& v_{\min} \leq v_k \leq v_{\max} \\
& a_{\min} \leq a_k \leq a_{\max} \\
& |\omega_k| \leq \omega_{\max} \\
& \|\mathbf{p}_k - \mathbf{c}^o\| \geq r_{\text{eff}}^o, \quad \forall \text{obstacles } o
\end{align}

\subsection{Cost Function}

\textbf{Stage Cost:}
\[
\ell(\mathbf{x}_k, \mathbf{u}_k, \mathbf{x}_k^{\text{tgt}}) = Q_p \|\mathbf{p}_k - \mathbf{p}_k^{\text{tgt}}\|^2 + Q_\theta (\theta_k - \theta_k^{\text{tgt}})^2 + R_a a_k^2 + R_\omega \omega_k^2
\]

\textbf{Terminal Cost:}
\[
\ell_f(\mathbf{x}_N, \mathbf{x}_N^{\text{tgt}}) = 10 Q_p \|\mathbf{p}_N - \mathbf{p}_N^{\text{tgt}}\|^2 + 10 Q_\theta (\theta_N - \theta_N^{\text{tgt}})^2
\]

\textbf{Weights (Tuned for Linear Paths):}
\begin{align}
Q_p &= 140.0 \\
Q_\theta &= 25.0 \quad \text{(increased for straight paths)} \\
R_a &= 0.01 \\
R_\omega &= 0.15 \quad \text{(increased to reduce curvature)} \\
Q_{\text{obs}} &= 900.0
\end{align}

\subsection{Uncertainty-Aware Constraints}

\textbf{Obstacle Inflation:}
\[
r_{\text{eff}}^o = r^o + k_\sigma \sqrt{\lambda_{\max}(\mathbf{P}^{\text{seek}})}
\]

where $k_\sigma = 0.5$ and $\lambda_{\max}(\mathbf{P}^{\text{seek}})$ is the largest eigenvalue of the seeker pose covariance.

\textbf{Velocity Scaling:}
\[
v_{\max}(t) = v_{\max}^{\text{base}} \cdot e^{-\alpha \sqrt{\lambda_{\max}(\mathbf{P}^{\text{seek}})}}
\]

with $\alpha = 0.5$.

\section{Simulation Implementation}

\subsection{Standalone Simulation Architecture}

The standalone simulation (`standalone_sim.py` and `animated_sim.py`) provides a complete testing environment without ROS2 dependencies.

\subsubsection{Map Generation}

\textbf{MapGenerator Class:}
\begin{itemize}
    \item Grid size: $100 \times 100$ cells
    \item Resolution: 0.05 m (5 cm)
    \item World bounds: $[-2.5, 2.5]$ m in both dimensions
    \item Ground truth map: Unknown to robot at startup
\end{itemize}

\textbf{Obstacle Placement:}
Strategic placement along the diagonal path from robot $(0, 0)$ to target $(1.5, 1.5)$:
\begin{itemize}
    \item Obstacle 1: $(0.5, 0.25)$ - blocks direct path
    \item Obstacle 2: $(0.25, 0.75)$ - blocks upper-left route
    \item Obstacle 3: $(0.9, 0.9)$ - blocks center of diagonal
    \item Obstacle 4: $(1.25, 0.5)$ - blocks lower-right route
    \item Obstacle 5: $(0.6, 1.4)$ - near target, blocks upper approach
\end{itemize}

\subsubsection{Fake LIDAR}

\textbf{FakeLidar Class:}
\begin{itemize}
    \item Number of beams: 360
    \item Max range: 3.5 m
    \item Min range: 0.12 m
    \item Range noise: $\sigma = 0.02$ m (Gaussian)
    \item Raycasting: Step size = resolution $\times$ 0.5
\end{itemize}

\textbf{Raycasting Algorithm:}
\begin{algorithm}[H]
\caption{Raycast in Occupancy Grid}
\begin{algorithmic}[1]
\Require Start position $(x, y)$, angle $\theta$, max range $r_{\max}$, map $\mathcal{M}$
\Ensure Range to first obstacle or $r_{\max}$
\State $d \gets 0$
\State $(x_c, y_c) \gets (x, y)$
\State $\Delta s \gets \text{resolution} \times 0.5$
\State $(\Delta x, \Delta y) \gets (\cos\theta \cdot \Delta s, \sin\theta \cdot \Delta s)$
\While{$d < r_{\max}$}
    \If{$\mathcal{M}[\text{grid}(x_c, y_c)]$ is occupied}
        \Return $d$
    \EndIf
    \State $(x_c, y_c) \gets (x_c + \Delta x, y_c + \Delta y)$
    \State $d \gets d + \Delta s$
\EndWhile
\Return $r_{\max}$
\end{algorithmic}
\end{algorithm}

\subsubsection{Robot Dynamics}

\textbf{Unicycle Model:}
\begin{align}
\dot{x} &= v \cos\theta \\
\dot{y} &= v \sin\theta \\
\dot{\theta} &= \omega
\end{align}

\textbf{Discrete-Time Integration:}
\begin{align}
x_{k+1} &= x_k + \Delta t \cdot v_k \cos\theta_k \\
y_{k+1} &= y_k + \Delta t \cdot v_k \sin\theta_k \\
\theta_{k+1} &= \theta_k + \Delta t \cdot \omega_k
\end{align}

\subsubsection{MCL Implementation (Simulation Only)}

\textbf{Note:} MCL is implemented for simulation and research but is \textbf{not used in hardware deployment}. The system uses EKF for robot localization in hardware.

\textbf{Particle Filter:}
\begin{itemize}
    \item Number of particles: 200-300
    \item Motion model: Unicycle with Gaussian noise
    \item Measurement model: LIDAR range likelihood
    \item Resampling: Systematic resampling
\end{itemize}

\textbf{Importance Weight:}
\[
w^{(i)} = \prod_{j=1}^{N_{\text{beams}}} \exp\left(-\frac{(r_j - \hat{r}_j^{(i)})^2}{2\sigma_r^2}\right)
\]

where $r_j$ is measured range and $\hat{r}_j^{(i)}$ is expected range for particle $i$.

\subsubsection{UKF for Target Tracking}

\textbf{State:} $[p_x, p_y, v_x, v_y]^T$

\textbf{Constant Velocity Model:}
\[
\mathbf{x}_{k+1} = \begin{bmatrix}
1 & 0 & \Delta t & 0 \\
0 & 1 & 0 & \Delta t \\
0 & 0 & 1 & 0 \\
0 & 0 & 0 & 1
\end{bmatrix} \mathbf{x}_k + \mathbf{w}_k
\]

\textbf{UKF Parameters:}
\begin{itemize}
    \item $\alpha = 0.001$ (spread of sigma points)
    \item $\beta = 2.0$ (incorporate prior knowledge)
    \item $\kappa = 0.0$ (secondary scaling parameter)
\end{itemize}

\subsection{ROS2 Simulation}

The ROS2 simulation (`simulation.launch.py`) provides ROS2 message-based testing.

\textbf{Components:}
\begin{itemize}
    \item \texttt{VoxelGridMapNode}: Creates test map, publishes \texttt{/map}
    \item \texttt{FakeLidarNode}: Simulates LIDAR, publishes \texttt{/scan}
    \item \texttt{MCLNode}: Localization (simulation only, not used in hardware)
    \item \texttt{UKFNode}: Target tracking, publishes \texttt{/target\_estimate}
    \item \texttt{MPCNode}: Control, publishes \texttt{/cmd\_vel}
\end{itemize}

\section{Hardware Deployment}

\subsection{Sensor Integration}

\subsubsection{LIDAR}

\textbf{Driver:} \texttt{rplidar\_ros} package
\begin{itemize}
    \item Topic: \texttt{/scan} (sensor\_msgs/LaserScan)
    \item Update rate: 10 Hz
    \item Range: 0.12 - 3.5 m
    \item Angular resolution: 1$^\circ$
\end{itemize}

\subsubsection{Camera}

\textbf{Driver:} \texttt{v4l2\_camera} package
\begin{itemize}
    \item Topic: \texttt{/camera/image\_raw} (sensor\_msgs/Image)
    \item Resolution: 640 $\times$ 480
    \item Frame rate: 30 Hz
    \item Color space: RGB
\end{itemize}

\subsubsection{IMU and Encoders}

\textbf{Topics:}
\begin{itemize}
    \item \texttt{/imu}: sensor\_msgs/Imu
    \item \texttt{/magnetic\_field}: sensor\_msgs/MagneticField
    \item \texttt{/odom}: nav\_msgs/Odometry (wheel encoders)
\end{itemize}

\subsection{Launch Sequence}

\textbf{Step 1: Start Cartographer}
\begin{lstlisting}[style=bash]
ros2 launch turtlebot3_cartographer cartographer.launch.py
\end{lstlisting}

\textbf{Step 2: Start Navigation System}
\begin{lstlisting}[style=bash]
ros2 launch turtlebot_interceptor single_robot_navigation.launch.py \
    goal_x:=1.5 goal_y:=1.5
\end{lstlisting}

\section{Performance Metrics}

\subsection{State Estimation}
\begin{itemize}
    \item EKF update rate: 100 Hz
    \item Position accuracy: $\pm$1 cm
    \item Heading accuracy: $\pm$2$^\circ$ (with magnetometer)
\end{itemize}

\subsection{Mapping}
\begin{itemize}
    \item Map update rate: 1-2 Hz
    \item Resolution: 2 cm
    \item Loop closure: Enabled
\end{itemize}

\subsection{Control}
\begin{itemize}
    \item MPC update rate: 10 Hz
    \item Horizon: 15 steps (1.5 s)
    \item Solver time: 20-50 ms (typically)
    \item Max solver time: $<$ 100 ms
\end{itemize}

\subsection{Target Tracking}
\begin{itemize}
    \item Camera detection range: $<$ 2.0 m
    \item LIDAR detection range: Up to 3.6 m
    \item Filter latency: $<$ 50 ms
    \item State propagation: Continuous when out of frame
\end{itemize}

\section{Challenges and Solutions}

\subsection{MPC Solver Failures}

\textbf{Problem:} CVXPY "Strict inequalities" error

\textbf{Solution:} Replaced conditional obstacle costs with smooth penalty functions:
\[
\text{cost} = \frac{w_{\text{obs}}}{d^2 + \epsilon}
\]

\subsection{Noisy LIDAR Data}

\textbf{Problem:} Hardware LIDAR produces noisy scans

\textbf{Solution:} Multi-layer filtering:
\begin{itemize}
    \item Exponential smoothing ($\alpha = 0.8$)
    \item Median filtering (last 10 detections)
    \item RANSAC-like outlier rejection (60\% max outliers)
\end{itemize}

\subsection{Camera Depth Estimation}

\textbf{Problem:} Target appearing closer than reality

\textbf{Solution:} Separate cone areas:
\begin{itemize}
    \item Yellow cones: 0.0208 m$^2$
    \item Blue target: 0.146 m$^2$ (7$\times$ larger)
\end{itemize}

\subsection{MPC Over-Curving}

\textbf{Problem:} Unnecessary curved paths

\textbf{Solution:} Increased angular penalty:
\begin{itemize}
    \item $R_\omega$: 0.01 $\rightarrow$ 0.15 (15$\times$ increase)
    \item $Q_\theta$: 10.0 $\rightarrow$ 25.0 (2.5$\times$ increase)
\end{itemize}

\section{Conclusion}

This system successfully demonstrates robust autonomous interception in unknown environments through:
\begin{itemize}
    \item \textbf{EKF-based robot localization} (100 Hz, multi-sensor fusion)
    \item Dual estimation architecture (EKF + Cartographer)
    \item Fast local grid for real-time obstacle awareness
    \item Multi-sensor target tracking (Vision + LIDAR)
    \item Uncertainty-aware MPC with adaptive constraints
    \item Robust filtering and state propagation
    \item Comprehensive error handling and fallbacks
\end{itemize}

The system has been validated from standalone simulation through ROS2 simulation to full hardware deployment, demonstrating production-ready robustness. \textbf{In hardware, EKF is used exclusively for robot localization} (MCL is available for simulation but not deployed).

\end{document}

